\section{Policy Response and Implementation}
\label{sec:policy_response}

\subsection{Current Policy Landscape}
According to the UNEP Adaptation Gap Report 2024:

\begin{itemize}
    \item 171 countries have national adaptation policies
    \item 51\% have secondary adaptation strategies
    \item 20\% have tertiary adaptation frameworks
\end{itemize}

In Africa specifically:
\begin{itemize}
    \item 45 countries have submitted updated Nationally Determined Contributions (NDCs)
    \item 38 countries have National Adaptation Plans (NAPs)
    \item 32 countries have integrated climate-health considerations into national policies
\end{itemize}

\subsection{Financial Mechanisms and Resources}
The African Union Climate Finance Report 2024 indicates:

\begin{itemize}
    \item \textbf{Current Funding:}
    \begin{itemize}
        \item Total climate finance flows to Africa: \$30 billion annually
        \item Only 10\% specifically allocated to health adaptation
        \item Significant gap between needs and available resources
    \end{itemize}
    
    \item \textbf{Funding Requirements:}
    \begin{itemize}
        \item Estimated \$38 trillion needed globally by 2050 for climate damage prevention
        \item African adaptation needs estimated at \$50 billion annually
        \item Health system resilience requires \$15 billion additional funding
    \end{itemize}
\end{itemize}

\subsection{Regional Implementation Progress}
African Ministers' 2024 Climate Resilience Strategy highlights:

\begin{itemize}
    \item \textbf{Policy Integration:}
    \begin{itemize}
        \item 75\% of countries have climate-health working groups
        \item 62\% have early warning systems in development
        \item 48\% have climate-resilient health facility programs
    \end{itemize}
    
    \item \textbf{Capacity Building:}
    \begin{itemize}
        \item Training programs for 50,000 healthcare workers
        \item Climate-health research initiatives in 28 countries
        \item Regional disease surveillance networks established
    \end{itemize}
\end{itemize}

\begin{keymessage}
\textbf{Policy Implementation Challenges:}
\begin{itemize}
    \item Significant funding gaps persist
    \item Limited technical capacity in many regions
    \item Need for better coordination between sectors
    \item Inadequate monitoring and evaluation systems
\end{itemize}
\end{keymessage}

\subsection{International Cooperation}
COP29 Africa Day (2024) outcomes include:

\begin{itemize}
    \item Enhanced commitment to climate finance
    \item Strengthened South-South cooperation
    \item New partnerships for technology transfer
    \item Regional climate-health research initiatives
\end{itemize}

\subsection{Policy Gaps and Barriers}
Key challenges identified in the 2024 AU Climate Strategy:

\begin{itemize}
    \item \textbf{Implementation Gaps:}
    \begin{itemize}
        \item Limited institutional capacity
        \item Insufficient technical resources
        \item Weak cross-sector coordination
    \end{itemize}
    
    \item \textbf{Resource Constraints:}
    \begin{itemize}
        \item Funding shortfalls
        \item Human resource limitations
        \item Technology access barriers
    \end{itemize}
\end{itemize}

\begin{recommendation}
\textbf{Priority Policy Actions:}
\begin{enumerate}
    \item Scale up climate finance mobilization
    \item Strengthen institutional capacity
    \item Enhance regional cooperation mechanisms
    \item Improve monitoring and evaluation systems
    \item Integrate health in all climate policies
\end{enumerate}
\end{recommendation}

\subsection{Future Policy Directions}
The African Union Climate Change Strategy 2025-2030 proposes:

\begin{itemize}
    \item Creation of an African Climate Health Observatory
    \item Development of regional early warning systems
    \item Establishment of climate-resilient health standards
    \item Enhanced capacity building programs
    \item Increased research and innovation funding
\end{itemize}

\subsection{Monitoring and Evaluation Framework}
Current tracking mechanisms include:

\begin{itemize}
    \item Annual progress reports on NDC implementation
    \item Health system resilience assessments
    \item Climate-health vulnerability indices
    \item Economic impact evaluations
\end{itemize}
